\section{The minimal cubic genus-one mechanism}

\subsection{Low-genus pullback capacity}

Let $\pi:S\to C$ be a relatively minimal non-isotrivial elliptic surface with
section, and put
\[
\chi=\chi(\mathcal O_S),\qquad g=g(C).
\]
The standard surface identities give
\[
e(S)=12\chi,\qquad q(S)=g,\qquad p_g(S)=\chi+g-1.
\]
Consequently
\[
b_2(S)=12\chi-2+4g,
\qquad
\boxed{h^{1,1}(S)=10\chi+2g.}
\]
If $R$ is the total root rank of reducible fibres, Shioda--Tate yields
\[
\boxed{
\rank E(\overline{\Q}(C))\le10\chi+2g-2-R.
}
\]
After a degree-$d$ base change of a rational elliptic surface,
$\chi=d$.  Hence
\[
\rank E(\overline{\Q}(C))\le10d+2g(C)-2-R.
\]

For a base with infinitely many rational points, the smallest possibilities
capable of rank thirty are
\[
\boxed{(d,g)=(4,0)\quad\text{or}\quad(d,g)=(3,1).}
\]
The second possibility is the decisive correction to the earlier assumption
that degree four was minimal.

\subsection{Ramification must avoid the discriminant}

Suppose the rational elliptic surface has twelve $I_1$ fibres.  Ramification
of index $e$ over an $I_1$ fibre produces an $I_e$ fibre and consumes root
rank $e-1$.  A degree-three genus-one cover has total ramification six by
Riemann--Hurwitz.  A cyclic cubic cover realizes this through three totally
ramified points.

Therefore a cubic pullback can reach the ceiling thirty only when all three
branch points lie over smooth fibres.  Branching over even one $I_1$ fibre
lowers the Shioda--Tate ceiling below thirty.

\subsection{An explicit positive-rank cyclic base}

Consider
\[
C_{36}:\qquad u^3=36t(t-1).
\]
The map $(t,u)\mapsto t$ is a cyclic cubic cover after adjoining a primitive
cube root of unity, totally ramified over $0,1,\infty$.  It is a genus-one
curve.  More generally, for
\[
u^3=c\,t(t-1),
\]
put
\[
X=4cu,\qquad Y=4c^2(2t-1).
\]
Then
\[
Y^2=X^3+16c^4.
\]
For $c=36$, the rational point $(t,u)=(3,6)$ maps to
\[
P=(864,25920).
\]
Doubling gives
\[
x(2P)=\frac{3456}{25},
\]
which is nonintegral.  If $P$ were torsion then $2P$ would be torsion, in
contradiction with Lutz--Nagell on the integral Mordell model.  Hence
\[
\boxed{\rank C_{36}(\Q)\ge1.}
\]
A M\"obius transformation moves the three branch values to any chosen triple
of rational smooth fibres.

\subsection{Cyclic cubic eigenspaces}

Let $C\to\mathbf P^1$ be such a cyclic cubic cover and let $\sigma$ generate
its deck group.  Over $\Q(\zeta_3)$, the Mordell--Weil space of the pullback
decomposes as
\[
M=M_1\oplus M_{\zeta}\oplus M_{\zeta^2}.
\]
The invariant part is the pullback of the original $E_8$ lattice, so
\[
\dim M_1=8.
\]
Complex conjugation exchanges the two nontrivial character spaces.  Writing
$\dim M_{\zeta}=\dim M_{\zeta^2}=r$, one obtains
\[
\boxed{\rank M=8+2r.}
\]
The Hodge bound gives $r\le11$.  In the absence of reducible fibres,
\[
\boxed{
\rank M=30
\iff r=11
\iff \rho(S_C)=32.
}
\]
The arithmetic problem is not merely Picard maximality: the thirty geometric
directions must descend to $\Q(C)$.
